%% Quantification along arbitrary set maps, including empty fibres.

\subsection{Powersets: substitution and quantifiers along a function}
\label{sec:p10-powerset-quantifiers}

Work in the category \(\mathsf{Set}\) of sets in the fixed universe
of Section~\ref{sec:p10-categorical-toolkit}.
All sets, maps, and indexed families below use that convention.
No nonemptiness hypothesis is imposed on these sets.

\begin{hbdefinition}{def:p10-powerset-quantifiers}{The powerset fibre and quantifiers}
For a set \(X\), let \(P(X)=\mathcal P(X)\), ordered by inclusion.
Its Boolean operations are intersection, union, complement in \(X\),
bottom \(\varnothing\), and top \(X\).
Indexed intersections are taken in the indicated powerset:
the empty intersection in \(P(X)\) is \(X\), and the empty
union is \(\varnothing\).
Its Heyting implication is
\[
                     U\Rightarrow_X V=(X\setminus U)\cup V.
\]
When \(X=\varnothing\), this is the one-element Boolean algebra.

For a function \(f:X\to Y\), define the reindexing map
\[
 f^*:P(Y)\longrightarrow P(X),\qquad
 f^*(V)=f^{-1}(V)=\{x\in X:f(x)\in V\}.
\]
Define two maps in the other direction by
\[
\begin{aligned}
 \exists_f(U)&=\{y\in Y:\exists x\in U\ (f(x)=y)\},\\
 \forall_f(U)&=\{y\in Y:\forall x\in X\
                                  (f(x)=y\Longrightarrow x\in U)\}.
\end{aligned}
\]
Thus \(\exists_f(U)\) is the image \(f[U]\), whereas
\(\forall_f(U)\) consists of points whose entire \(f\)-fibre
lies in \(U\).
\index{powerset!quantifiers}\index{inverse image!as substitution}
\index{quantifier!along a function}
\end{hbdefinition}

\begin{hbtheorem}{thm:p10-powerset-adjunctions}{proved=annex:p10.quant-set.1}{The two powerset adjunctions and their function laws}
Inverse image preserves the Boolean operations and all unions
and intersections.  The maps \(\exists_f,f^*,\forall_f\) are
monotone, and
\[
                      \exists_f\dashv f^*\dashv\forall_f.
\]
Explicitly, for \(U\subseteq X\) and \(V\subseteq Y\),
\[
 \exists_f(U)\subseteq V\ \Longleftrightarrow\ U\subseteq f^*(V),
 \qquad
 f^*(V)\subseteq U\ \Longleftrightarrow\ V\subseteq\forall_f(U).
\]
These properties determine both quantifier maps uniquely.
For \(f:X\to Y\), \(g:Y\to Z\), the function laws are
\[
\begin{gathered}
 (1_X)^*=\exists_{1_X}=\forall_{1_X}=1_{P(X)},\\
 (g\circ f)^*=f^*\circ g^*,\qquad
 \exists_{g\circ f}=\exists_g\circ\exists_f,\qquad
 \forall_{g\circ f}=\forall_g\circ\forall_f .
\end{gathered}
\]
Moreover,
\[
 \forall_f(U)=Y\setminus\exists_f(X\setminus U),\qquad
 \exists_f\Bigl(\bigcup_j U_j\Bigr)=\bigcup_j\exists_f(U_j),\qquad
 \forall_f\Bigl(\bigcap_j U_j\Bigr)=\bigcap_j\forall_f(U_j).
\]
The family may be empty.  In particular,
\[
 \exists_f(\varnothing)=\varnothing,\quad
 \forall_f(X)=Y,\quad
 \exists_f(X)=f[X],\quad
 \forall_f(\varnothing)=Y\setminus f[X].
\]
\end{hbtheorem}

\begin{hbexample}{def:p10-powerset-fibres}{A universal quantifier includes an empty fibre}
Let \(X=\{0,1,2\}\), \(Y=\{a,b,c\}\), and
\(f(0)=f(1)=a\), \(f(2)=b\).
The fibres at \(a,b,c\) are respectively \(\{0,1\},\{2\},\varnothing\).
For \(U=\{0,2\}\),
\[
                    \exists_f(U)=\{a,b\},\qquad
                    \forall_f(U)=\{b,c\}.
\]
The fibre at \(a\) meets \(U\) but is not contained in it;
the empty fibre at \(c\) is contained in \(U\).
For a projection \(\pi:I\times A\to I\), the same definitions give
\[
 \exists_\pi(U)=\{i:\exists a\in A\ ((i,a)\in U)\},\qquad
 \forall_\pi(U)=\{i:\forall a\in A\ ((i,a)\in U)\}.
\]
If \(A=\varnothing\), its only domain predicate is
\(\varnothing\), with
\(\exists_\pi(\varnothing)=\varnothing\) and
\(\forall_\pi(\varnothing)=I\).
\end{hbexample}

The powerset and inverse-image construction is
\cite[Example~5.2.2, p.~42]{Pitts1995Draft};
the projection quantifiers are Example~5.5.3, pp.~53--54.
Lawvere's powerset example uses inverse image for substitution
and direct image for existential quantification
\cite[p.~13]{Lawvere2006AdjointnessReprint}.
These locators refer respectively to the 18 May 1995 draft
and the authorized 2006 reprint.

\subsection{Pullback stability and conjunction}
\label{sec:p10-powerset-base-change}

\begin{hbtheorem}{thm:p10-powerset-base-change}{proved=annex:p10.quant-set.2}{Beck--Chevalley and Frobenius for powersets}
Suppose the following is a pullback square of sets:
\[
\begin{array}{ccc}
 X'&\xrightarrow{\ g'\ }&X\\
 {\scriptstyle f'}\downarrow&&\downarrow{\scriptstyle f}\\
 Y'&\xrightarrow{\ g\ }&Y.
\end{array}
\]
Then both Beck--Chevalley identities hold, as maps
\(P(X)\to P(Y')\):
\[
              g^*\circ\exists_f=\exists_{f'}\circ(g')^*,
              \qquad
              g^*\circ\forall_f=\forall_{f'}\circ(g')^*.
\]
For an arbitrary function \(f:X\to Y\), Frobenius reciprocity is
\[
             \exists_f\bigl(U\cap f^*(V)\bigr)
                         =\exists_f(U)\cap V
                \quad(U\subseteq X,\ V\subseteq Y).
\]
For a square which is only assumed to commute, the following
inclusions always hold for \(U\subseteq X\):
\[
 \exists_{f'}((g')^*U)\subseteq g^*(\exists_fU),
 \qquad
 g^*(\forall_fU)\subseteq\forall_{f'}((g')^*U).
\]
Commutativity alone does not imply the reverse inclusions.
\end{hbtheorem}

\begin{hbexample}{def:p10-powerset-commuting-square}{A commuting square without base change}
In the displayed shape, take
\(X'=\varnothing\), \(X=Y=Y'=\{*\}\), \(f=g=1_{\{*\}}\),
and let \(f',g'\) be the unique maps from the empty set.
The square commutes, but is not a pullback: the compatible
pair of points of \(X,Y'\) has no point over it in \(X'\).
For \(U=X\),
\[
             g^*(\exists_fU)=\{*\},\qquad
             \exists_{f'}((g')^*U)=\varnothing.
\]
For \(U=\varnothing\),
\[
             g^*(\forall_fU)=\varnothing,\qquad
             \forall_{f'}((g')^*U)=\{*\}.
\]
Thus each automatic inclusion can be strict.
\end{hbexample}

The general-square stability equations and the automatic
comparison directions are discussed in
\cite[Remark~5.6.6, pp.~58--59]{Pitts1995Draft};
Frobenius is equation~(5.10) and Remark~5.5.2, p.~53.
The proof below establishes both identities for every pullback
of set maps.

\subsection{Diagonal equality and the Boolean doctrine}
\label{sec:p10-powerset-equality}

\begin{hbtheorem}{thm:p10-powerset-equality}{proved=annex:p10.quant-set.3}{Powerset equality with parameters}
For a set \(A\), let \(\Delta_A:A\to A\times A\) send
\(a\) to \((a,a)\), and put
\[
              \delta_A=\exists_{\Delta_A}(A)
                         =\{(a,a):a\in A\}\in P(A\times A).
\]
For any \(s,t:I\to A\),
\[
             \langle s,t\rangle^*\delta_A
                        =\{i\in I:s(i)=t(i)\}.
\]
In particular \(\Delta_A^*\delta_A=A\).

Use the left-associated product \((I\times A)\times A\)
and the maps
\[
\begin{aligned}
 e:I\times A&\longrightarrow(I\times A)\times A,
                         &e(i,a)&=((i,a),a),\\
 p:(I\times A)\times A&\longrightarrow I\times A,
                         &p((i,a),b)&=(i,a),\\
 r:(I\times A)\times A&\longrightarrow A\times A,
                         &r((i,a),b)&=(a,b).
\end{aligned}
\]
For \(U\subseteq I\times A\),
\[
                    \exists_e(U)=p^*(U)\cap r^*(\delta_A).
\]
Consequently, for \(V\subseteq(I\times A)\times A\),
\[
       p^*(U)\cap r^*(\delta_A)\subseteq V
                   \quad\Longleftrightarrow\quad U\subseteq e^*(V).
\]
The powerset functor with these predicates is therefore a
Boolean first-order hyperdoctrine with equality, in the precise
sense of Definitions~\ref{def:p10-first-order-hyperdoctrine}
and~\ref{def:p10-doctrine-equality}.
It satisfies the parameter-product Beck--Chevalley identities
and the full parameterized diagonal adjunction, including
empty parameter sets and empty sorts.
\end{hbtheorem}

Pitts gives the parameterized equality condition and the Set
diagonal in
\cite[\S5.6, equation~(5.12), Definition~5.6.1, and
Example~5.6.3, pp.~55--56]{Pitts1995Draft}.
Lawvere describes equality through existential quantification
along the diagonal in
\cite[p.~12]{Lawvere2006AdjointnessReprint}.

The terminal set interpreting a context with no variables is
a singleton \(1\).  It is distinct from an empty sort \(A=\varnothing\).
For such a sort, quantification along \(A\to1\) sends its top
predicate \(A=\varnothing\) to \(\varnothing\) existentially
and to \(1\) universally.
The doctrine above permits these empty sorts; interpretations
of ordinary first-order logic with nonempty domains impose
nonemptiness as an additional condition.
